\documentclass{plantillaucsd}
\usepackage[utf8]{inputenc}
\usepackage[T1]{fontenc}
\usepackage{geometry}
\usepackage{graphicx}
\usepackage{float}
\usepackage{hyperref}
\usepackage[english,indonesian]{babel}

\title{Optimasi Pemrosesan Template LaTeX Menggunakan Pendekatan Berbasis AI}
\author{Academic Support Agent}
\date{\today}

\begin{document}
\maketitle

\begin{abstract}

\end{abstract}

\chapter{PENDAHULUAN}


\chapter{TINJAUAN PUSTAKA}


% Original content (commented out):
% 
\documentclass{plantillaucsd}

% ************************* Declaraciones ***********************

% Poner aquí el título de la tesis
\title{Evaluación de Lorem ipsum dolor sit amet, consectetuer adipiscing elit ut purus elit}

% Toda esta información es necesaria. 
\author{Nombre Apellido} % Nombre entero del autor
\apellido{Apellido} % Apellido, para numerar las páginas
\date{Mes, Año} % Fecha, en formato {Mes, Año}
\matricula{20xx-0001} 

\facultad{Facultad de [Llenar]}
\escuela{Escuela de [Llenar]}
\titulo{Licenciatura en [Llenar]} % El título académico que se recibirá.

% ******************* Inicio del Documento *********************

\begin{document}

\titlepage{}

\chapter*{Agradecimientos}
\thispagestyle{empty}
Aquí van los agradecimientos.
\newpage
\chapter*{Dedicatorias}
Aquí van las dedicatorias.
\thispagestyle{empty}
\newpage


\tableofcontents{}
\thispagestyle{empty}
\newpage

\input{capitulos/introduccion}
\input{capitulos/capitulo1}
\input{capitulos/capitulo2}
\input{capitulos/capitulo3}
\input{capitulos/conclusion}

\nocite{*}

\printbibliography[heading=bibintoc]

\end{document}

\end{document}